\chapter{Dos gestos y una ventana}

\inicio{B}{enito} Arce era un hombre corpulento con manos de estibador y voz de niño de coro, y aquella tarde estaba convencido de que iba a pasar el resto de su vida en una mina de hielo.

---Yo cambié el cartel ---repetía---\nobreak. Lo cambié como lo he cambiado diez mil veces. Grabé el cristal nuevo, sumé en su hoja, cambié el número del cartel y resté en la hoja del viejo. Como dice el manual. Como dice la puerta.

---Le creo ---dijo Inés---\nobreak. Por eso le he hecho venir. Necesito que lo haga otra vez, despacio, delante de mí.

Estaban en el Gran Salón, a medianoche, que era la hora en que menos lectores había. El cartel de la carta de aproximación al muelle doce brillaba en la pared, a la altura de los ojos. Era una placa rectangular dividida en sesenta y cuatro casillas, cada una de las cuales podía estar encendida o apagada. Las cuarenta y ocho de la derecha formaban el número del nicho. Las dieciséis de la izquierda estaban apagadas.

---¿Siempre han estado así? ---preguntó Inés, señalándolas.

---Siempre, doctora ---dijo Siete---\nobreak. Los nichos tienen cuarenta y ocho cifras. Los carteles tienen sesenta y cuatro casillas porque los fundadores los fabricaron en serie, todos iguales. Las dieciséis primeras nunca se han encendido.

Inés se quedó mirando las casillas apagadas un momento más de lo necesario. Luego se volvió hacia Beto.

---Adelante.

\escena
Lo representaron como una obra de teatro. Siete hizo de ANA-3; Beto hizo de sí mismo; Inés se sentó en un banco con una libreta y el comisionado Mendieta, que había insistido en acompañarlos, se quedó de pie con los brazos cruzados.

Para que se viera bien, usaron un cristal de pruebas en el nicho más cercano, a cuarenta metros del cartel, y Siete se comprometió a andar despacio.

---Momento cero ---dictó Inés---\nobreak. El cartel anuncia el cristal A. La hoja de A marca uno: el uno del cartel. Nadie más lo usa. Siete, lee el cartel.

Siete se acercó a la placa, la miró y dijo en voz alta:

---Nicho A.

Y echó a andar hacia el pasillo.

---Momento uno ---dijo Inés---\nobreak. Beto, ahora.

Beto ya tenía preparado el cristal B. Lo colocó en su nicho, sumó uno en su hoja, se acercó al cartel y, con un único movimiento de la palanca, cambió el número: el cartel pasó a anunciar el nicho B. Luego fue al nicho A y restó uno en su hoja.

---Cero ---anunció.

---Y en la vida real ---dijo Inés---\nobreak, ahora el cristal A bajaría al fundidor. Supongamos que lo ha hecho. Siete, sigue.

Siete llegó al nicho A, alargó la mano hacia la hoja y se detuvo. Sus ojos ámbar parpadearon.

---Doctora, si el cristal ya no está, la hoja tampoco. Si la toco, toco la hoja de lo que haya ahora en el nicho.

---Exacto. Eso es lo que hizo ANA-3.

Inés pasó la página de la libreta y escribió la escena como una tabla. Siempre pensaba mejor en tablas.

\begin{pizarra}
\begin{tabular}{@{}c p{2.6cm} p{2.6cm} c c@{}}
\toprule
& \textbf{Siete (lector)} & \textbf{Beto (archivero)} & \textbf{Cartel} & \textbf{Hoja de A} \\
\midrule
0 & & & A & 1 \\
1 & Lee el cartel: «A» & & A & 1 \\
2 & (camina) & Apunta 1 en B, pone B en el cartel & B & 1 \\
3 & (camina) & Resta 1 en A: queda 0. ¡Al fundidor! & B & --- \\
4 & Suma 1 en la hoja de A \ \textcolor{oxido}{$\boldsymbol\times$} & & B & fundido \\
\bottomrule
\end{tabular}
\end{pizarra}

---Ahí está toda la historia ---dijo---\nobreak. El lector necesita dos gestos: leer el cartel y apuntarse en la hoja. Entre el primero y el segundo hay una ventana. Si el archivero se cuela por esa ventana, el cristal se funde con un lector en camino.

---Pues que el lector se apunte primero ---dijo Mendieta.

---¿En qué hoja, comisionado? Para saber qué hoja tiene que firmar, primero tiene que leer el cartel.

---Pues que lea el cartel y firme a la vez.

Inés le dirigió la mirada que reservaba para los estudiantes que, sin saberlo, hacían la pregunta correcta.

---Esa es la pregunta. Siete, explícale al comisionado por qué no se puede.

\escena
---Los fundadores ---explicó Siete--- construyeron cada cartel y cada hoja con un mecanismo especial. Ese mecanismo sabe hacer ciertos gestos \emph{de golpe}: nadie puede ver un estado intermedio ni colarse en medio. Si dos personas tiran a la vez de la misma palanca, el mecanismo atiende primero a una y luego a la otra, en algún orden, pero nunca mezcla los dos gestos.

Fue enumerándolos con los dedos, que en su caso eran de aleación de titanio.

---El primero es \emph{arrancar}: sumar o restar una cantidad a lo que marca la placa, y ver qué marcaba antes, todo en un solo movimiento. El segundo es \emph{trocar}: poner un valor nuevo y quedarse con el viejo en la mano. Beto lo usa para cambiar el cartel. El tercero es \emph{cotejar}: cambiar el valor solo si sigue siendo el que uno esperaba, y si no, ver qué hay.

---¿Y entonces?

---Entonces cada gesto actúa sobre una placa, señor. Sobre un cartel, o sobre una hoja. El cartel está en el Gran Salón; la hoja, en la Bóveda, a veces a tres kilómetros. No existe ningún gesto que lea un cartel y modifique una hoja a la vez. Haría falta una mano con dos brazos que llegaran a la vez a dos sitios distintos.

---Y esa mano ---añadió Inés--- no la tiene nadie. Ni aquí ni en ningún otro sitio.

Mendieta se rascó la nuca.

---Doctora, ¿con qué frecuencia pasa esto? Nueve milésimas de segundo me parece muy poca ventana para colar un archivero.

---Muy poca. Es rarísimo que ocurra en una lectura concreta. Pero el Archivo atiende millones de lecturas al día y cambia miles de carteles. Haga la multiplicación. Siete, ¿cuántos incidentes del fundidor tiene registrados el Archivo como «avería»?

Siete permaneció inmóvil un instante.

---Cuatro mil ciento doce en los últimos diez años, doctora. Cristales fundidos dos veces, hojas que marcan más usuarios de los que hay, lecturas con datos «corruptos». Todos clasificados como fallos mecánicos.

---Ninguno lo era ---dijo Inés---\nobreak. Todos eran esto. Un fallo que solo aparece cuando dos gestos se intercalan de cierta manera es el peor tipo de fallo que existe: pasa todas las pruebas, porque en las pruebas casi nunca coinciden, y luego aparece en el peor momento, lejos de su causa. A veces es un dato equivocado. A veces es un fundidor atascado. Y a veces es una nave a nueve kilómetros de un muelle.

\escena
Beto Arce, que llevaba un rato sentado en el suelo con la espalda contra la pared, levantó la cabeza.

---Entonces no voy a la mina.

---No va a la mina ---dijo Inés---\nobreak. Lo que usted hizo estaba bien. Lo que está mal es el Reglamento. O, mejor dicho, la forma en que lo cumplimos.

---Ya se ha intentado arreglar otras veces ---dijo Siete---\nobreak. Los registros mencionan varias propuestas. Ninguna se adoptó.

---Me lo imagino. ---Inés se puso en pie y se sacudió la tiza de las rodillas---\nobreak. Siempre hay tres arreglos obvios, y los tres rompen una Regla distinta. Hagamos la prueba. Comisionado, ¿qué haría usted?

Mendieta lo pensó.

---Que el lector compruebe, cuando llega al nicho, si el cristal sigue siendo el suyo.

---¿Cómo lo compruebas, Siete?

---Leyéndolo, doctora. Y si ya no es el suyo, al leerlo ya he leído lo que no debía. La comprobación misma es el error.

---Primer arreglo descartado. ¿Otro?

---Que el fundidor espere un poco antes de fundir ---dijo Beto desde el suelo---\nobreak. Un minuto, por ejemplo. Así da tiempo a que lleguen los rezagados.

---Funciona ---admitió Inés---\nobreak, si sabe cuánto esperar. Pero entonces la Bóveda se llena de cristales que nadie lee, esperando su turno. Rompe la Tercera Regla. Volveremos a esto, porque es más interesante de lo que parece. ¿Siete?

---Que el cartel tenga una llave ---dijo el robot---\nobreak. Quien quiera leerlo o cambiarlo, debe tener la llave. Así nadie cambia el cartel mientras otro está a mitad de camino.

---Funciona. Y si el que tiene la llave se detiene ---porque se le acaba la batería, porque le llama otra tarea, porque se le olvida---\nobreak, todos los demás esperan. Para siempre, si hace falta. Rompe la Segunda Regla.

Siete inclinó la cabeza, asimilando.

---Entonces cada arreglo cumple dos Reglas y rompe una.

---Así es. Y el sistema que tenemos ahora cumple la Segunda y la Tercera, y rompe la Primera de vez en cuando, sin avisar. ---Inés miró la placa del cartel, con sus dieciséis casillas apagadas a la izquierda---\nobreak. Mañana la Directora convocará al Consejo, y el Consejo propondrá la llave, porque es lo que siempre se propone. Y yo tengo hasta entonces para encontrar algo mejor.

---¿Y si no lo encuentra? ---preguntó Mendieta.

---Entonces pondrán la llave, y dentro de un año usted volverá a mi despacho con otra nave. Pero esa vez la nave no se habrá estrellado. Simplemente no habrá llegado nunca, porque se habrá quedado esperando la llave.

Se marchó sin despedirse. Siete se quedó un momento más mirando el cartel, imitando sin saberlo la postura de la doctora. Luego, porque era un robot concienzudo, contó las casillas apagadas. Seguían siendo dieciséis.
